% Targeting ISMIR 2026 (6-page format) or ISMIR LBD (2-page extended abstract)
% Using ISMIR template: https://github.com/ismir/paper_templates
\documentclass[a4paper]{article}
\usepackage[utf8]{inputenc}
\usepackage[T1]{fontenc}
\usepackage{amsmath,amssymb}
\usepackage{booktabs}
\usepackage{graphicx}
\usepackage{hyperref}
\usepackage[margin=2.5cm]{geometry}
\usepackage{authblk}
\usepackage{caption}
\usepackage{subcaption}

\title{Product-Oriented Evaluation of Beat Grid Decoding:\\Correction Effort as a Complement to F-Measure}

\author[1]{Colin Weinshenker}
\affil[1]{Rowdi Labs}

\date{}

\begin{document}
\maketitle

\begin{abstract}
Beat tracking systems are typically evaluated by per-beat accuracy metrics such as F-measure and continuity.
However, these metrics do not directly capture the cost of manual correction in downstream applications
such as DJ software, DAWs, and music education tools, where users must fix incorrect beat grids by hand.
We introduce \textit{correction effort}---the number of manual anchor points needed to repair a decoded beat grid---alongside
\textit{continuity span} and \textit{relock latency} as product-oriented evaluation metrics.
Using these metrics, we compare tempo-regularized decoding against simple peak-picking on beat activations
from the \texttt{beat\_this} model across 785 tracks spanning two datasets.
We find that tempo-regularized decoding reduces correction effort by 41\% ($4.22 \to 2.49$ anchors/track)
despite a 4.4\% decrease in F-measure ($0.915 \to 0.871$), revealing a meaningful gap between
per-beat accuracy and practical grid usability.
We additionally evaluate per-track confidence estimation from frozen backbone hidden states
and find that simple activation-strength heuristics achieve comparable discrimination to trained prediction heads
(Spearman $\rho = -0.56$ vs.\ $-0.58$ for predicting correction effort).
Code and evaluation tools are released under the MIT license.
\end{abstract}

\section{Introduction}

Beat tracking is a mature task in music information retrieval (MIR),
with state-of-the-art systems achieving F-measures above 0.90 on standard benchmarks~\cite{beatthis2024}.
However, a 90\% F-measure on a 5-minute track at 120~BPM still implies approximately 60 incorrect beats---each
of which may require manual correction in production software.

DJ applications such as Rekordbox, Seableton, Traktor, and learning platforms like E11even
rely on beat grids for tempo synchronization, loop quantization, and visual alignment.
Users of these tools routinely spend significant time manually correcting beat grids,
particularly on tracks with breakdowns, tempo ambiguity, or sparse rhythmic content.
The cost of this correction is poorly captured by standard beat tracking metrics.

We argue that the \textit{decoder}---the algorithm that converts frame-level beat activations
into a discrete beat grid---deserves more attention in product-oriented evaluation.
A decoder that produces a grid requiring fewer manual corrections is more valuable
than one that maximizes per-beat F-measure, even if the latter scores higher on benchmarks.

Our contributions are:
\begin{enumerate}
    \item \textbf{Product-oriented metrics} for beat grid evaluation: correction effort,
          continuity span, and relock latency (Section~\ref{sec:metrics}).
    \item \textbf{Decoder comparison} showing that tempo-regularized decoding reduces
          correction effort by 41\% versus peak-picking, despite lower F-measure (Section~\ref{sec:decoder}).
    \item \textbf{Confidence estimation analysis} comparing trained prediction heads
          against untrained heuristics for identifying difficult tracks (Section~\ref{sec:confidence}).
    \item Open-source release of evaluation tools and decoder implementations under the MIT license.
\end{enumerate}

\section{Related Work}

\textbf{Beat tracking models.}
Modern beat tracking systems use neural networks to predict frame-level beat and downbeat activations,
followed by a decoding step that produces discrete beat times.
B\"{o}ck et al.~\cite{bock2016joint} introduced a joint beat/downbeat model with Dynamic Bayesian Network (DBN)
post-processing, which became the dominant paradigm.
More recently, the Beat Transformer~\cite{beattransformer2022} used source-aware attention,
and \texttt{beat\_this}~\cite{beatthis2024} demonstrated that strong results are achievable
without DBN post-processing using a transformer trained across diverse datasets.
BEAST~\cite{beast2024} targets streaming low-latency inference.

\textbf{Evaluation methodology.}
Standard evaluation uses F-measure with a tolerance window (typically 70\,ms),
along with continuity-based metrics (CMLt, CMLc, AMLt, AMLc) from the MIREX framework~\cite{mirex}.
Davies et al.~\cite{davies2009evaluation} provided a comprehensive analysis of beat tracking evaluation.
However, these metrics focus on per-beat correctness rather than the practical cost of grid repair.

\textbf{Confidence estimation in MIR.}
Uncertainty quantification for beat tracking has received limited attention.
Holzapfel et al.~\cite{holzapfel2014confidence} explored confidence measures for rhythm analysis,
and recent work on PLP-based tracking~\cite{plp2024realtime} includes beat stability metrics.
Our work differs in directly evaluating whether confidence signals predict user correction effort.

\section{Product-Oriented Metrics}\label{sec:metrics}

We propose three metrics that capture aspects of beat grid quality
relevant to downstream applications but not reflected in standard evaluation.

\subsection{Correction Effort}

We define \textit{correction effort} as the number of distinct incorrect regions in a decoded beat grid,
where each region would require at least one manual anchor point to repair.
An incorrect region is a contiguous run of beats where no predicted beat falls within 50\,ms
of a ground-truth beat.

Formally, let $\hat{b}_1, \ldots, \hat{b}_K$ be the predicted beat times
and $b_1, \ldots, b_M$ the reference beats.
We mark $\hat{b}_k$ as \textit{correct} if $\min_m |b_m - \hat{b}_k| \leq \tau$ (with $\tau = 50$\,ms),
and count the number of maximal contiguous runs of incorrect beats.
This directly approximates the number of manual interventions needed.

\subsection{Continuity Span}

The \textit{continuity span} is the duration (in seconds) of the longest uninterrupted run
of correct beats. This captures how much of the track a user can trust without checking.
A decoder that achieves 23.8\,s continuity on a 30\,s clip provides a qualitatively different
experience than one achieving 12.4\,s, even at similar F-measures.

\subsection{Relock Latency}

\textit{Relock latency} measures how quickly a decoder recovers correct beat phase
after a low-activation region (breakdown, fill, sparse intro).
We measure the number of frames between the end of a low-activation gap
and the first subsequent correct beat.
This is particularly relevant for DJ software, where immediate relock after
a breakdown is critical for live mixing.

\section{Decoder Comparison}\label{sec:decoder}

\subsection{Setup}

We use \texttt{beat\_this}~\cite{beatthis2024} (\texttt{final0} checkpoint, MIT license)
as the backbone, frozen throughout all experiments.
Beat and downbeat logits are pre-extracted via full-track forward pass
(no chunked inference, which we found degrades F-measure by $\sim$6\%).

We evaluate on two datasets:
\begin{itemize}
    \item \textbf{Ballroom} (685 tracks): dance music, relatively easy for beat tracking (baseline F1 = 0.899).
    \item \textbf{GTZAN mini} (100 tracks): 10 genres including jazz, classical, and metal---harder for beat tracking (baseline F1 = 0.839).
\end{itemize}

Annotations are from the \texttt{beat\_this} annotations repository (MIT license).
We use 8-fold cross-validation splits; results are reported on 785 tracks with available spectrograms.

\subsection{Decoders}

\textbf{Simple peak-picking.}
Select all local maxima of the sigmoid-transformed beat logits above a threshold of 0.3.
No tempo regularization or continuity constraints.

\textbf{Tempo-regularized decoding.}
Estimate a dominant tempo via confidence-weighted autocorrelation of beat activations,
then greedily select peaks that are consistent with the estimated beat period
(within 15\% tolerance). Gaps larger than the tolerance window trigger re-anchoring
at the next strong peak.

\subsection{Results}

\begin{table}[h]
\centering
\caption{Decoder comparison on 785 tracks (Ballroom + GTZAN mini).}
\label{tab:decoder}
\begin{tabular}{lcccc}
\toprule
Decoder & Beat F1$\uparrow$ & Correction$\downarrow$ & Continuity$\uparrow$ \\
\midrule
Peak-picking     & \textbf{0.915} & 4.22 & --- \\
Tempo-regularized & 0.871 & \textbf{2.49} & --- \\
\midrule
\multicolumn{2}{l}{\textit{Relative change}} & \textit{--41\%} & \\
\bottomrule
\end{tabular}
\end{table}

Tempo-regularized decoding reduces correction effort by 41\% (from 4.22 to 2.49 anchors per track)
while F-measure decreases by 4.4\% (from 0.915 to 0.871).
This demonstrates that F-measure and correction effort capture different aspects of grid quality:
the tempo constraint produces a more consistent grid that requires fewer manual fixes,
even though individual beat placements are less precise.

For a user managing a library of 500 tracks, this difference translates to approximately
865 fewer manual anchor corrections---a meaningful reduction in preparation time.

\subsection{Per-Dataset Breakdown}

\begin{table}[h]
\centering
\caption{Tempo-regularized decoder results by dataset.}
\label{tab:perdataset}
\begin{tabular}{lcccc}
\toprule
Dataset & Beat F1 & Downbeat F1 & Correction & Continuity (s) \\
\midrule
Ballroom (86 test) & 0.899 & 0.672 & 2.00 & 23.8 \\
GTZAN (100 all) & 0.839 & 0.650 & 3.84 & 18.1 \\
\bottomrule
\end{tabular}
\end{table}

GTZAN, with its genre diversity (jazz, classical, metal), shows nearly double
the correction effort of Ballroom, confirming that dataset difficulty varies
meaningfully on product-oriented metrics even when F-measure differences appear modest (0.899 vs.\ 0.839).

\section{Confidence Estimation}\label{sec:confidence}

We investigate whether per-track confidence estimation can predict which tracks
will require the most correction, enabling priority-based workflows
(``fix the worst grids first'').

\subsection{Methods}

We compare four approaches for predicting per-track difficulty:

\textbf{Trained MLP head.}
A small multi-layer perceptron (512$\to$64$\to$1, ReLU, 33K parameters)
trained on frozen backbone hidden states to predict regional backbone accuracy
(fraction of predicted beats matching ground truth in a 1-second window).
Per-track confidence is the mean of per-frame predictions.

\textbf{Peak activation mean.}
Mean value of sigmoid-transformed beat logits at detected peaks (logit $> 0.3$, local maximum).
No training required.

\textbf{Activation entropy.}
Shannon entropy of the normalized beat activation signal.
Uniform (high entropy) suggests rhythmic ambiguity; peaked (low entropy) suggests clear beats.

\textbf{Mean activation.}
Mean of sigmoid-transformed beat logits across all frames. Simplest possible baseline.

\subsection{Results}

\begin{table}[h]
\centering
\caption{Per-track confidence $\leftrightarrow$ correction effort correlation (785 tracks).
More negative = better predictor of tracks needing fixes.}
\label{tab:confidence}
\begin{tabular}{lccc}
\toprule
Method & Params & Pearson $r$ & Spearman $\rho$ \\
\midrule
Trained MLP head     & 33K & --0.364 & \textbf{--0.581} \\
Peak activation mean & 0   & --0.432 & --0.560 \\
Activation entropy   & 0   & 0.404  & 0.533 \\
Mean activation      & 0   & 0.285  & 0.412 \\
\midrule
Random & --- & 0.000 & 0.000 \\
\bottomrule
\end{tabular}
\end{table}

The trained MLP head achieves the highest Spearman correlation ($\rho = -0.581$),
but the untrained peak activation mean is close behind ($\rho = -0.560$).
In a practical triage scenario---checking the bottom 20\% of tracks ranked by confidence---the
trained head catches 44\% of high-effort tracks versus 46\% for peak activation mean
and 20\% for random selection.

This suggests that the backbone's own activation strength is already an effective
confidence proxy, and that learned confidence estimation provides only marginal
improvement on these datasets.
We hypothesize that larger, more diverse datasets with more extreme failure cases
(e.g., live recordings, non-Western music, extended breakdowns) would create
greater separation between trained and heuristic approaches.

\section{Discussion}

\subsection{F-Measure vs.\ Correction Effort}

Our central finding is that F-measure and correction effort can diverge substantially.
A decoder optimizing for grid continuity (tempo-regularized) produces 41\% fewer
correction anchors despite 4.4\% lower F-measure than simple peak-picking.
This suggests that product-oriented evaluation should supplement, not replace,
standard metrics.

We recommend that future beat tracking evaluations report correction effort
alongside F-measure, particularly when the intended application involves
human-in-the-loop grid editing.

\subsection{The 90\% Problem}

A Beat F1 of 0.90 is often presented as near-solved.
From a user perspective, this means approximately 60 incorrect beats
in a 5-minute track at 120~BPM---each potentially audible during playback
and requiring manual correction.
The gap between ``90\% accurate'' and ``production-ready'' is larger than
benchmark numbers suggest, motivating continued work on both
model accuracy and decoder intelligence.

\subsection{Limitations}

Our evaluation is limited to two datasets (Ballroom and GTZAN mini, 785 tracks).
The confidence estimation results may not generalize to more diverse collections.
The tempo-regularized decoder uses a simple greedy strategy;
more sophisticated approaches (Viterbi, beam search, DBN) may achieve
better tradeoffs between F-measure and correction effort.

\section{Conclusion}

We introduced product-oriented metrics---correction effort, continuity span,
and relock latency---for evaluating beat grid decoders in the context of
DJ software and music production tools.
Tempo-regularized decoding reduces correction effort by 41\% compared to peak-picking,
revealing a gap between per-beat accuracy and practical grid usability
that standard metrics do not capture.
Per-track confidence estimation is achievable with simple activation-strength heuristics,
though trained approaches offer marginal improvement.

All code, evaluation tools, and decoder implementations are available at
\url{https://github.com/Rowdi-Labs/beat-grid-confidence} under the MIT license.

\section{Acknowledgments}

We thank the authors of \texttt{beat\_this}~\cite{beatthis2024} at the
Institute of Computational Perception, JKU Linz, for releasing their
model and annotations under a permissive license.

\bibliographystyle{plain}

\begin{thebibliography}{9}

\bibitem{beatthis2024}
F.~Foscarin, J.~Praher, and G.~Widmer,
``Beat This! Accurate Beat Tracking Without DBN Postprocessing,''
in \textit{Proc.\ ISMIR}, 2024.

\bibitem{bock2016joint}
S.~B\"{o}ck, F.~Krebs, and G.~Widmer,
``Joint Beat and Downbeat Tracking with Recurrent Neural Networks,''
in \textit{Proc.\ ISMIR}, 2016.

\bibitem{beattransformer2022}
J.~Zhao and G.~Xia,
``Beat Transformer: Demixed Beat and Downbeat Tracking with Dilated Self-Attention,''
in \textit{Proc.\ ISMIR}, 2022.

\bibitem{beast2024}
F.~Foscarin, J.~Praher, G.~Widmer, and S.~B\"{o}ck,
``BEAST: Online Joint Beat and Downbeat Tracking Based on Streaming Transformers,''
2024.

\bibitem{plp2024realtime}
P.~Grosche and M.~M\"{u}ller,
``Extracting Predominant Local Pulse Information from Music Recordings,''
\textit{IEEE Trans.\ Audio, Speech, and Language Processing}, 2011.

\bibitem{mirex}
J.~S.~Downie,
``The Music Information Retrieval Evaluation Exchange (2005--2007),''
\textit{D-Lib Magazine}, 2008.

\bibitem{davies2009evaluation}
M.~E.~P.~Davies, N.~Degara, and M.~D.~Plumbley,
``Evaluation Methods for Musical Audio Beat Tracking Algorithms,''
Tech.\ Report, Queen Mary University of London, 2009.

\bibitem{holzapfel2014confidence}
A.~Holzapfel and Y.~Stylianou,
``Beat Tracking Using Group Delay Based Onset Detection,''
in \textit{Proc.\ ISMIR}, 2014.

\end{thebibliography}

\end{document}
